\textbf{Phần 1. Trắc nghiệm lựa chọn}
\setcounterpageref{ex}{0}

\begin{ex} %[5D12N69]
Kết quả của phép tính $45{,}6 + 18{,}75$ là
\choice
{$63{,}35$}
{\True $64{,}35$}
{$64{,}45$}
{$63{,}45$}
\loigiai{
$45{,}6 + 18{,}75 = 64{,}35$.\\
\textbf{Đáp án của câu hỏi là B}
}
\end{ex}

\begin{ex} %[5D12N69]
Kết quả của phép tính $\dfrac{2}{3} + \dfrac{1}{6}$ là
\choice
{$\dfrac{3}{9}$}
{$\dfrac{4}{6}$}
{\True $\dfrac{5}{6}$}
{$1$}
\loigiai{
$\dfrac{2}{3} + \dfrac{1}{6} = \dfrac{4}{6} + \dfrac{1}{6} = \dfrac{5}{6}$.\\
\textbf{Đáp án của câu hỏi là C}
}
\end{ex}

\begin{ex} %[5D12H69]
Kết quả của phép nhân $3{,}25 \times 4$ là
\choice
{$12{,}5$}
{\True $13$}
{$13{,}25$}
{$12$}
\loigiai{
$3{,}25 \times 4 = 13$.\\
\textbf{Đáp án của câu hỏi là B}
}
\end{ex}

\begin{ex} %[5D12H69]
Kết quả của phép tính $18{,}6 : 0{,}1$ là
\choice
{$1{,}86$}
{$0{,}186$}
{\True $186$}
{$18{,}6$}
\loigiai{
Chia cho $0{,}1$ tức là nhân với $10$: $18{,}6 \times 10 = 186$.\\
\textbf{Đáp án của câu hỏi là C}
}
\end{ex}

\begin{ex} %[5D12H69]
Tính nhanh giá trị biểu thức: $3{,}6 \times 45 + 3{,}6 \times 55$.
\choice
{$36$}
{$3,600$}
{$180$}
{\True $360$}
\loigiai{
$3{,}6 \times (45 + 55) = 3{,}6 \times 100 = 360$.\\
\textbf{Đáp án của câu hỏi là D}
}
\end{ex}

\begin{ex} %[5D12H69]
Kết quả của phép chia $\dfrac{4}{5} : \dfrac{2}{3}$ là
\choice
{\True $\dfrac{6}{5}$}
{$\dfrac{8}{15}$}
{$\dfrac{5}{6}$}
{$\dfrac{2}{5}$}
\loigiai{
$\dfrac{4}{5} \times \dfrac{3}{2} = \dfrac{12}{10} = \dfrac{6}{5}$.\\
\textbf{Đáp án của câu hỏi là A}
}
\end{ex}

\begin{ex} %[5D12V69]
Tìm $x$ biết: $x \times 2{,}5 = 48{,}5 - 16$.
\choice
{$12$}
{$14$}
{$15$}
{\True $13$}
\loigiai{
$x \times 2{,}5 = 32{,}5 \Rightarrow x = 32{,}5 : 2{,}5 = 13$.\\
\textbf{Đáp án của câu hỏi là D}
}
\end{ex}

\begin{ex} %[5D12V69]
Giá trị của biểu thức $(12{,}4 - 4{,}9) \times 3{,}2 + 6$ là
\choice
{\True $30$}
{$28$}
{$32$}
{$34$}
\loigiai{
$7{,}5 \times 3{,}2 + 6 = 24 + 6 = 30$.\\
\textbf{Đáp án của câu hỏi là A}
}
\end{ex}

\begin{ex} %[5D12V69]
Một mảnh đất hình chữ nhật có chiều dài $24{,}5\text{ m}$, chiều rộng $12\text{ m}$. Diện tích mảnh đất đó là
\choice
{$284\text{ m}^2$}
{$304\text{ m}^2$}
{\True $294\text{ m}^2$}
{$290\text{ m}^2$}
\loigiai{
$24{,}5 \times 12 = 294\text{ (m}^2\text{)}$.\\
\textbf{Đáp án của câu hỏi là C}
}
\end{ex}

\begin{ex} %[5D12V69]
Tính giá trị của biểu thức: $\dfrac{3}{4} \times \dfrac{8}{9} + \dfrac{1}{3}$.
\choice
{\True $1$}
{$\dfrac{2}{3}$}
{$\dfrac{4}{3}$}
{$\dfrac{5}{6}$}
\loigiai{
$\dfrac{2}{3} + \dfrac{1}{3} = 1$.\\
\textbf{Đáp án của câu hỏi là A}
}
\end{ex}

\begin{ex} %[5D12V69]
Một người mua $3{,}5\text{ kg}$ đường loại $24\,000$ đồng/kg và $2\text{ kg}$ bột mì loại $18\,000$ đồng/kg. Người đó đưa cho người bán hàng tờ $200\,000$ đồng. Số tiền thừa nhận lại là
\choice
{$70\,000$ đồng}
{$90\,000$ đồng}
{$85\,000$ đồng}
{\True $80\,000$ đồng}
\loigiai{
Tổng tiền mua: $3{,}5 \times 24\,000 + 2 \times 18\,000 = 84\,000 + 36\,000 = 120\,000$ (đồng).\\
Tiền thừa: $200\,000 - 120\,000 = 80\,000$ (đồng).\\
\textbf{Đáp án của câu hỏi là D}
}
\end{ex}

\begin{ex} %[5D12C69]
Tính giá trị biểu thức bằng cách thuận tiện: $\dfrac{0{,}25 \times 61{,}4 \times 40}{0{,}5 \times 0{,}2 \times 100}$. Kết quả là
\choice
{$6{,}14$}
{\True $61{,}4$}
{$614$}
{$0{,}614$}
\loigiai{
Tử số: $(0{,}25 \times 40) \times 61{,}4 = 10 \times 61{,}4 = 614$.\\
Mẫu số: $(0{,}5 \times 0{,}2) \times 100 = 0{,}1 \times 100 = 10$.\\
Kết quả: $614 : 10 = 61{,}4$.\\
\textbf{Đáp án của câu hỏi là B}
}
\end{ex}

\textbf{Phần 2. Trắc nghiệm đúng sai}
\setcounterpageref{ex}{0}

\begin{ex} %[5D12H69]
Xét tính đúng sai của các phép tính sau:
\choiceTF
{\True $15{,}4 \times 100 = 1\,540$.}
{$45{,}8 : 0{,}01 = 4{,}58$.}
{\True $\dfrac{5}{7} \times \dfrac{14}{15} = \dfrac{2}{3}$.}
{\True $0{,}8 : 0{,}25 = 3{,}2$.}
\loigiai{
\begin{itemize}
    \item $15{,}4 \times 100 = 1540$. Đúng.
    \item $45{,}8 : 0{,}01 = 4580 \neq 4{,}58$. Sai.
    \item $\dfrac{5 \times 14}{7 \times 15} = \dfrac{2}{3}$. Đúng.
    \item $0{,}8 : 0{,}25 = 3{,}2$. Đúng.
\end{itemize}
\begin{itemize}
    \item a) \textbf{Đúng}.
    \item b) \textbf{Sai}.
    \item c) \textbf{Đúng}.
    \item d) \textbf{Đúng}.
\end{itemize}
}
\end{ex}

\begin{ex} %[5D12V69]
Cho biểu thức $A = 4{,}8 \times 3{,}5 + 4{,}8 \times 6{,}5$ và $B = (12{,}6 - 2{,}6) \times 4{,}8$. Xét tính đúng sai:
\choiceTF
{\True Giá trị biểu thức $A = 48$.}
{\True Giá trị biểu thức $B = 48$.}
{Giá trị của biểu thức $A$ lớn hơn giá trị của biểu thức $B$.}
{\True $A = B$.}
\loigiai{
\begin{itemize}
    \item $A = 4{,}8 \times (3{,}5 + 6{,}5) = 4{,}8 \times 10 = 48$. Đúng.
    \item $B = 10 \times 4{,}8 = 48$. Đúng.
    \item $A = B = 48$ nên khẳng định $A > B$ là Sai.
    \item $A = B$. Đúng.
\end{itemize}
\begin{itemize}
    \item a) \textbf{Đúng}.
    \item b) \textbf{Đúng}.
    \item c) \textbf{Sai}.
    \item d) \textbf{Đúng}.
\end{itemize}
}
\end{ex}

\textbf{Phần 3. Trả lời ngắn}
\setcounterpageref{ex}{0}

\begin{ex} %[5D12H69]
Tính kết quả phép tính: $45{,}8 + 34{,}75$.
\shortans{80,55}
\loigiai{
\textbf{Đáp số:} $80{,}55$\\
$45{,}8 + 34{,}75 = 80{,}55$.
}
\end{ex}

\begin{ex} %[5D12H69]
Tính kết quả: $12{,}5 \times 0{,}8$.
\shortans{10}
\loigiai{
\textbf{Đáp số:} $10$\\
$12{,}5 \times 0{,}8 = 10$.
}
\end{ex}

\begin{ex} %[5D12V69]
Tính giá trị biểu thức: $150 - 24{,}5 \times 4$.
\shortans{52}
\loigiai{
\textbf{Đáp số:} $52$\\
$150 - 98 = 52$.
}
\end{ex}

\begin{ex} %[5D12V69]
Tìm $x$ biết: $x : 3{,}6 = 12{,}5$.
\shortans{45}
\loigiai{
\textbf{Đáp số:} $45$\\
$x = 12{,}5 \times 3{,}6 = 45$.
}
\end{ex}

\begin{ex} %[5D12V69]
Tính nhanh: $2{,}5 \times 17{,}8 \times 4$.
\shortans{178}
\loigiai{
\textbf{Đáp số:} $178$\\
$(2{,}5 \times 4) \times 17{,}8 = 10 \times 17{,}8 = 178$.
}
\end{ex}

\begin{ex} %[5D12C69]
Một kho hàng có $120$ tấn gạo. Ngày thứ nhất xuất $\dfrac{1}{3}$ số gạo, ngày thứ hai xuất $40\%$ số gạo còn lại. Hỏi sau hai ngày kho hàng còn lại bao nhiêu tấn gạo?
\shortans{48}
\loigiai{
\textbf{Đáp số:} $48$\\
Ngày 1 xuất: $120 : 3 = 40$ tấn. Còn lại: $120 - 40 = 80$ tấn.\\
Ngày 2 xuất: $80 \times 40\% = 32$ tấn.\\
Còn lại sau 2 ngày: $80 - 32 = 48$ tấn.
}
\end{ex}
